% Generic template for validation runs of SMOS L2 against a gridded reference
% (e.g. ERA5-Land, SMAP L3). SMOS L2 is dataset 0, the reference dataset 1 in
% the netcdf results. All names, file names and labels are taken from the run's
% ContentVars, so the same template can be used for multiple runs.
\def\runidx{\detokenize{$<ContentVars['report_run_index']>$}}
\def\runrefname{\detokenize{$<ContentVars['NetcdfMetaVars']['val_dc_dataset_pretty_name1']>$}}
\def\runpair{0-\detokenize{$<ContentVars['NetcdfMetaVars']['val_dc_dataset0']>$}_and_1-\detokenize{$<ContentVars['NetcdfMetaVars']['val_dc_dataset1']>$}}

\subsection{Validation run \runidx{} -- \runrefname{} (\detokenize{$<ContentVars['NetcdfMetaVars']['val_dc_version_pretty_name1']>$})}

\subsubsection{Validation settings}

The following validation settings were used for the comparison of \ac{SMOS} \ac{L2} against \runrefname.
For a full list of the settings used, see the results page of this validation run, linked in Table \ref{tab:val_run_list}.

\begin{itemize}
    \item Spatial reference (spatial matching): \textbf{\detokenize{$<ContentVars['ConfigVars']['spatial_reference']>$}}
    \item Temporal reference (temporal matching): \textbf{\detokenize{$<ContentVars['ConfigVars']['temporal_reference']>$}}
    \item Temporal matching window: \textbf{\detokenize{$<ContentVars['ConfigVars']['temporal_matching']>$} hours}
    \item Scaling reference: \textbf{\detokenize{$<ContentVars['ConfigVars']['scaling_reference']>$}}
    \item Validation metrics \detokenize{$<'\\textbf{are}' if ContentVars['ConfigVars']['anomalies_method']!='none' else '\\textbf{are not}'>$} based on anomalies.
\end{itemize}

\subsubsection{Results}

\noindent \ac{SMOS} \ac{L2} was compared against \runrefname{} at \textbf{\detokenize{$<ContentVars['NetcdfVars']['n_gpis']>$}} grid points.
The datasets are spatially matched via a nearest-neighbour search, with \detokenize{$<ContentVars['ConfigVars']['spatial_reference']>$} as the spatial reference; results are therefore computed on the \detokenize{$<ContentVars['ConfigVars']['spatial_reference']>$} grid.
Validation metrics could be computed for \textbf{\detokenize{$<ContentVars['NetcdfVars']['n_ok'][0]>$}} grid points (\detokenize{$<ContentVars['NetcdfVars']['percent_ok'][0]>$}\,\%).
For the remaining \detokenize{$<ContentVars['NetcdfVars']['n_failed'][0]>$} grid points, no metrics were computed, mostly because too few valid observations remained after filtering and temporal matching (Fig. \ref{fig:status\runidx}).

\vspace{3mm}

\noindent Table \ref{tab:sumstats\runidx} lists summary statistics across all valid grid points: the number of temporally matched observations used for the metric computation ("\#Obs."), Pearson's R, the bias (difference in means) and the \ac{ubRMSD}.
Additional metrics are available in the validation results \href{\detokenize{$<ContentVars['NetcdfMetaVars']['url']>$}}{online}.

\begin{table}[htb!]
    \centering
    \caption{Summary statistics of \ac{SMOS} \ac{L2} against \runrefname{} across all valid grid points.}
    \label{tab:sumstats\runidx}
    \begin{tabular}{ l r r r r }
        \toprule
        \textbf{Metric} & \textbf{\#Obs.} & \textbf{Pearson's R} & \textbf{Bias}  & \textbf{ubRMSD}  \\
        \midrule
        \textbf{Mean}   & \detokenize{$<ContentVars['SummaryVars']['MEAN_N. OBS.']>$}   & \detokenize{$<ContentVars['SummaryVars']['MEAN_PEARSON'S R']>$}     & \detokenize{$<ContentVars['SummaryVars']['MEAN_BIAS']>$}   & \detokenize{$<ContentVars['SummaryVars']['MEAN_UNBIASED ROOT-MEAN-SQUARE DEVIATION']>$}   \\
        \textbf{Median}   & \detokenize{$<ContentVars['SummaryVars']['MEDIAN_N. OBS.']>$}   & \detokenize{$<ContentVars['SummaryVars']['MEDIAN_PEARSON'S R']>$}     & \detokenize{$<ContentVars['SummaryVars']['MEDIAN_BIAS']>$}   & \detokenize{$<ContentVars['SummaryVars']['MEDIAN_UNBIASED ROOT-MEAN-SQUARE DEVIATION']>$}   \\
        \textbf{Interquartile range}   & \detokenize{$<ContentVars['SummaryVars']['IQ RANGE_N. OBS.']>$}   & \detokenize{$<ContentVars['SummaryVars']['IQ RANGE_PEARSON'S R']>$}     & \detokenize{$<ContentVars['SummaryVars']['IQ RANGE_BIAS']>$}   & \detokenize{$<ContentVars['SummaryVars']['IQ RANGE_UNBIASED ROOT-MEAN-SQUARE DEVIATION']>$}   \\
        \bottomrule
    \end{tabular}
\end{table}

\begin{figure}[ht!]
\centering
\begin{minipage}[c]{0.34\textwidth}
    \centering
    \plotfig{./qa4sm_graphics/bulk_barplot_status_\runpair.png}
    \subcaption{} % (a)
\end{minipage}%
\hfill
\begin{minipage}[c]{0.64\textwidth}
    \centering
    \mapfig{./qa4sm_graphics/bulk_overview_\runpair_status.png}
    \subcaption{} % (b)
\end{minipage}
\caption{Number of successfully validated grid points and errors (a), and their location (b) in the validation run against \runrefname.}
\label{fig:status\runidx}
\end{figure}

\noindent Validation metrics are computed from temporally matched observations of \ac{SMOS} \ac{L2} and \runrefname.
A \detokenize{$<ContentVars['ConfigVars']['temporal_matching']>$}-hour temporal matching window is applied to each \detokenize{$<ContentVars['ConfigVars']['temporal_reference']>$} observation to find the closest data point in the other dataset.
The resulting sample size, which applies to all computed metrics, is shown in Fig. \ref{fig:nobs\runidx}.

\begin{figure}[ht!]
\centering
\begin{minipage}[c]{0.34\textwidth}
    \centering
    \plotfig{./qa4sm_graphics/bulk_boxplot_n_obs.png}
    \subcaption{} % (a)
\end{minipage}%
\hfill
\begin{minipage}[c]{0.64\textwidth}
    \centering
    \mapfig{./qa4sm_graphics/bulk_overview_n_obs.png}
    \subcaption{} % (b)
\end{minipage}
\caption{Distribution of the validation sample size (a) and sample size per location (b) in the validation run against \runrefname.}
\label{fig:nobs\runidx}
\end{figure}

\noindent Fig. \ref{fig:Rbox\runidx} shows Pearson's R between \ac{SMOS} \ac{L2} and \runrefname{} at all valid grid points.
The \ac{ubRMSD} is shown in Fig. \ref{fig:urmsd\runidx}.

\begin{figure}[ht!]
\centering
\begin{minipage}[c]{0.34\textwidth}
    \centering
    \plotfig{./qa4sm_graphics/bulk_boxplot_R.png}
    \subcaption{} % (a)
\end{minipage}%
\hfill
\begin{minipage}[c]{0.64\textwidth}
    \centering
    \mapfig{./qa4sm_graphics/bulk_overview_\runpair_R.png}
    \subcaption{} % (b)
\end{minipage}
\caption{Distribution of Pearson's R across all valid grid points (a) and spatial overview of Pearson's R (b) between \ac{SMOS} \ac{L2} and \runrefname.}
\label{fig:Rbox\runidx}
\end{figure}

\begin{figure}[ht!]
\centering
\begin{minipage}[c]{0.34\textwidth}
    \centering
    \plotfig{./qa4sm_graphics/bulk_boxplot_urmsd.png}
    \subcaption{} % (a)
\end{minipage}%
\hfill
\begin{minipage}[c]{0.64\textwidth}
    \centering
    \mapfig{./qa4sm_graphics/bulk_overview_\runpair_urmsd.png}
    \subcaption{} % (b)
\end{minipage}
\caption{Distribution of the \ac{ubRMSD} across all valid grid points (a) and spatial overview of the \ac{ubRMSD} (b) between \ac{SMOS} \ac{L2} and \runrefname.}
\label{fig:urmsd\runidx}
\end{figure}
